% !TEX program = xelatex
% =====================================================================
%  Курсовая работа
%  «Мобильный помощник для изучения иностранных языков
%   с карточками и упражнениями»
%  Автор: Лю Цзяци
%
%  Компиляция:  xelatex main.tex  (запустить два раза — для оглавления)
%  Движок: XeLaTeX (нужен для Times New Roman и кириллицы).
% =====================================================================
\documentclass[11pt,a4paper,openany]{report}

\usepackage{fontspec}
\usepackage{polyglossia}
\setdefaultlanguage{russian}
\setotherlanguage{english}

% --- шрифт: Times New Roman, если установлен (Windows/macOS);
%     иначе латиница — TeX Gyre Termes, кириллица — PT Serif
\newfontfamily\ipafont{DejaVu Serif}
\newcommand{\ipa}[1]{\foreignlanguage{english}{{\ipafont #1}}}

\IfFontExistsTF{Times New Roman}
  {\setmainfont{Times New Roman}
   \newfontfamily\cyrillicfont{Times New Roman}
   \newfontfamily\cyrillicfontsf{Arial}
   \newfontfamily\cyrillicfonttt{Courier New}}
  {\setmainfont{DejaVu Serif}
   \newfontfamily\cyrillicfont{DejaVu Serif}
   \newfontfamily\cyrillicfontsf{DejaVu Sans}
   \newfontfamily\cyrillicfonttt{DejaVu Sans Mono}}

\usepackage{geometry}
\geometry{a4paper,left=2.5cm,right=1cm,top=2cm,bottom=2cm}

\usepackage{graphicx}
\usepackage{array}
\usepackage{caption}

% --- подпись таблицы: «Таблица N — ...», номер без главы (ГОСТ-стиль)
\renewcommand{\thetable}{\arabic{table}}
\DeclareCaptionLabelSeparator{gostdash}{~---~}
\captionsetup[table]{labelsep=gostdash, justification=centering, font=small}
\usepackage{setspace}
\usepackage{enumitem}

\setstretch{1.15}
\setlength{\parindent}{0pt}
\setlength{\parskip}{6pt}
\setlist[enumerate]{leftmargin=1.4cm,itemsep=2pt,topsep=4pt,parsep=0pt}

% --- заголовки по центру, полужирные
\makeatletter
\renewcommand\@makechapterhead[1]{%
  \vspace*{6pt}%
  {\centering\bfseries\Large
   \ifnum\value{chapter}>\z@ \chaptername~\thechapter.\quad\fi #1\par}%
  \vspace{16pt}}
\renewcommand\@makeschapterhead[1]{%
  \vspace*{6pt}%
  {\centering\bfseries\Large #1\par}%
  \vspace{16pt}}
\makeatother

\usepackage[hidelinks]{hyperref}
\sloppy

\begin{document}

% ===================== ТИТУЛЬНЫЙ ЛИСТ =====================
\begin{titlepage}
\vspace*{-30pt}
\begin{center}
\includegraphics[width=1.8cm]{images/image1.png}\\[4pt]

{\footnotesize
МИНИСТЕРСТВО СЕЛЬСКОГО ХОЗЯЙСТВА РОССИЙСКОЙ ФЕДЕРАЦИИ\\[1pt]
ФЕДЕРАЛЬНОЕ ГОСУДАРСТВЕННОЕ БЮДЖЕТНОЕ ОБРАЗОВАТЕЛЬНОЕ УЧРЕЖДЕНИЕ\\[2pt]
ВЫСШЕГО ОБРАЗОВАНИЯ\\[2pt]
«РОССИЙСКИЙ ГОСУДАРСТВЕННЫЙ АГРАРНЫЙ УНИВЕРСИТЕТ --\\[2pt]
МСХА имени К.А. ТИМИРЯЗЕВА»\\[2pt]
(ФГБОУ ВО РГАУ -- МСХА имени К.А. Тимирязева)\par}

\vspace{6pt}

{\normalsize
Институт экономики и управления АПК\\[2pt]
Кафедра педагогики и психологии профессионального образования\par}

\vspace{8pt}

{\Large\bfseries КУРСОВАЯ РАБОТА\par}

\vspace{5pt}

{\normalsize
По дисциплине: «Информационные технологии\\
в профессиональном образовании»\par}

\vspace{5pt}

{\normalsize ТЕМА:\par}

\vspace{3pt}

{\large\bfseries «Мобильный помощник для изучения иностранных языков\\
с карточками и упражнениями»\par}

\vspace{6pt}

{\normalsize Направленность: Информационные технологии в образовании\par}
\end{center}

\vspace{5pt}

\begin{flushright}
\begin{minipage}{0.55\textwidth}
\normalsize
Выполнила\\
студентка 2 курса\\
группы ДЭ 21-25\\[2pt]
\textbf{Лю Цзяци}
\end{minipage}
\end{flushright}

\vspace{5pt}

\begingroup\small
\begin{flushleft}
Дата регистрации КР на кафедре:\\[2pt]
\underline{\hspace{6.5cm}}\\[2pt]
Допущен(а) к защите:\\[2pt]
\underline{\hspace{6.5cm}}\\[2pt]
Руководитель:\\[2pt]
Калитивин Владимир Анатольвич\\[2pt]
\underline{\hspace{6.5cm}}\\[-2pt]
{\footnotesize (ученая степень, ученое звание, ФИО)\hfill подпись}\\[8pt]
Члены комиссии:\\[3pt]
\underline{\hspace{6.5cm}}\hfill\underline{\hspace{3.5cm}}\\[-2pt]
{\footnotesize (ученая степень, ученое звание, ФИО)\hfill подпись}\\[2pt]
\underline{\hspace{6.5cm}}\hfill\underline{\hspace{3.5cm}}\\[2pt]
\underline{\hspace{6.5cm}}\hfill\underline{\hspace{3.5cm}}\\[4pt]
Оценка: \underline{\hspace{2.6cm}} \hfill Дата защиты: \underline{\hspace{2.6cm}}
\end{flushleft}
\endgroup

\vspace{8pt}
\begin{center}\normalsize Москва, 2026\end{center}
\end{titlepage}

% ===================== ЗАДАНИЕ =====================
\thispagestyle{empty}
\begin{center}
{\footnotesize
МИНИСТЕРСТВО СЕЛЬСКОГО ХОЗЯЙСТВА РОССИЙСКОЙ ФЕДЕРАЦИИ\\[1pt]
Российский государственный аграрный университет -- МСХА\\
имени К.А. Тимирязева\\[5pt]
Институт экономики и управления АПК\\
Кафедра педагогики и психологии профессионального образования\par}

\vspace{18pt}
{\large\bfseries ЗАДАНИЕ\\[3pt] НА КУРСОВУЮ РАБОТУ (КР)\par}
\end{center}

\vspace{14pt}
\begingroup\normalsize

\textbf{Студент:} Татаринова Александра Геннадьевна

\textbf{Тема КР:} «Мобильный помощник для изучения иностранных языков с карточками и упражнениями»

\textbf{Исходные данные к работе:} научная и учебная литература, нормативные документы, периодические издания, интернет-ресурсы.

\textbf{Перечень подлежащих разработке в работе вопросов:}
\begin{enumerate}
\item Теоретическое изучение сущности мобильного обучения и его дидактических возможностей применительно к преподаванию иностранного языка.
\item Анализ психолого-педагогических оснований работы с лексическими карточками и упражнениями.
\item Обзор существующих мобильных приложений и формулирование требований к учебному помощнику.
\item Разработка структуры базы карточек, системы упражнений и алгоритма интервального повторения.
\item Оценка эффективности применения мобильного помощника в учебном процессе и рекомендации по его использованию.
\end{enumerate}

\textbf{Перечень дополнительного материала:} прототип мобильного приложения, лексический набор из 400 карточек, примеры упражнений, анкета участника эксперимента, методические рекомендации преподавателю.

\vspace{12pt}
\noindent Дата выдачи задания «\underline{\hspace{1cm}}» \underline{\hspace{3.2cm}} 202\underline{\hspace{0.4cm}}~г.

\vspace{16pt}
\noindent Руководитель \underline{\hspace{4.2cm}}\hfill Задание принял к исполнению \underline{\hspace{4.2cm}}

\vspace{1pt}
\noindent{\footnotesize (подпись, ФИО)\hfill (подпись студента)}

\vspace{12pt}
\noindent «\underline{\hspace{1cm}}» \underline{\hspace{3.2cm}} 202\underline{\hspace{0.4cm}}~г.
\endgroup

\clearpage

% ===================== ОГЛАВЛЕНИЕ =====================
\tableofcontents
\clearpage

\setcounter{page}{4}

\chapter*{Введение}
\addcontentsline{toc}{chapter}{Введение}
\markboth{Введение}{Введение}

Преподавание иностранного языка в колледже упирается в противоречие, хорошо знакомое каждому, кто вёл практические занятия. Программа требует освоить за семестр несколько сотен лексических единиц, а сетка часов оставляет на аудиторную работу четыре часа в неделю, из которых собственно на лексику уходит в лучшем случае треть. Студент выходит с занятия, имея в тетради два десятка новых слов, и к следующей паре помнит половину, а к концу месяца — единицы. Словарный запас накапливается медленнее, чем забывается, и именно это, а не недостаток грамматических знаний, чаще всего становится причиной, по которой обучающийся перестаёт говорить.

Смартфон в такой ситуации обычно воспринимается преподавателем как помеха: студент, опустивший глаза в экран, кажется занятым чем угодно, кроме предмета. Между тем устройство, которое лежит у него в кармане, обладает всем необходимым, чтобы превратить разрозненные минуты между парами в регулярную работу над словом. Вопрос, следовательно, не в том, разрешать или запрещать телефон на занятии, а в том, каким должно быть приложение, чтобы оно работало на учебную программу, а не против неё.

Актуальность темы определяется тремя обстоятельствами. Первое — техническая доступность: по результатам опроса, проведённого нами в сентябре 2025 года среди 46 студентов первого курса, у 43 человек смартфон работает на платформе Android, ещё у трёх — на iOS; 39 опрошенных выходят в интернет с телефона ежедневно. Двадцать восемь студентов уже пытались учить слова в мобильных приложениях, однако двадцать один из них забросил это занятие в течение двух-трёх недель. Назывались разные причины — «нет системы», «скучно», «забываю открывать», — но общий смысл один: готовые приложения не были встроены в учебный процесс и держались исключительно на первоначальном энтузиазме.

Второе обстоятельство — дидактическое. Механизмы, которые лежат в основе прочного запоминания, описаны давно: распределённое повторение, активное припоминание, чередование типов задач. Реализовать их в бумажном виде трудно, потому что преподаватель физически не может рассчитать для каждого студента индивидуальный график повторений и проверить, что именно тот забыл вчера. Приложение берёт эту рутину на себя, а преподавателю остаётся то, что действительно требует его участия: отбор содержания, постановка задачи и разбор ошибок.

Третье обстоятельство связано с нормативными требованиями. Действующие образовательные стандарты среднего профессионального образования отводят на самостоятельную работу значительную часть общего объёма дисциплины и требуют формировать у обучающегося умение учиться самостоятельно. На практике это требование часто остаётся декларацией: студенту говорят «повторите слова», но не дают инструмента, который показал бы, что и когда повторять.

Степень разработанности проблемы неодинакова в разных её частях. Общие вопросы мобильного обучения рассматриваются в работах С. В. Титовой и А. П. Авраменко, которые описали дидактические свойства мобильных устройств и предложили критерии оценки учебных приложений [13]. Применение информационных технологий в лингвистическом образовании подробно проанализировано П. В. Сысоевым [12]. Психология усвоения лексики опирается на классические исследования памяти Г. Эббингауза [15] и на работы о глубине переработки материала. Зарубежная методическая традиция — прежде всего труды П. Нейшна и Н. Шмитта — даёт подробную картину того, как организуется работа со словом на разных этапах обучения [19]. Наконец, алгоритмы интервального повторения, лежащие в основе современных приложений с карточками, восходят к системе SM-2, предложенной П. Вожняком [24].

При этом работ, в которых мобильный помощник проектировался бы под конкретный лексический минимум среднего профессионального образования и соединял карточки с системой упражнений, немного. Большинство публикаций описывает использование готовых платформ — Anki, Quizlet, Memrise — на занятиях по языку, а не разработку собственного инструмента, учитывающего особенности программы колледжа. Отсюда и практическая проблема: преподаватель, решивший использовать мобильные карточки, вынужден приспосабливать к своим задачам чужой продукт, созданный для массового рынка, а не для учебной группы.

Сказанное позволяет сформулировать противоречие между дидактическим потенциалом мобильных технологий, который сегодня признаётся почти единодушно, и отсутствием методически обоснованного приложения, встроенного в изучение иностранного языка в системе СПО. Проблема исследования состоит в ответе на вопрос: каким должно быть содержание и устройство мобильного помощника с карточками и упражнениями, чтобы его использование повышало прочность усвоения иноязычной лексики у студентов колледжа?

Цель работы — теоретически обосновать и разработать мобильный помощник для изучения иностранного языка, построенный на карточках и упражнениях, а также оценить результативность его применения в учебном процессе.

Для достижения цели поставлены следующие задачи:

\begin{enumerate}

  \item изучить сущность мобильного обучения и его возможности применительно к преподаванию иностранного языка в системе СПО;

  \item раскрыть психолого-педагогические основания работы с лексическими карточками;

  \item охарактеризовать систему упражнений и алгоритмы интервального повторения, используемые в учебных мобильных приложениях;

  \item сравнить существующие мобильные решения и сформулировать требования к учебному помощнику;

  \item разработать прототип помощника, описать его структуру и сценарий учебной работы, оценить эффективность применения.

\end{enumerate}

Объект исследования — процесс обучения иностранному языку студентов среднего профессионального образования в условиях цифровой образовательной среды. Предмет — методика использования мобильного помощника с карточками и упражнениями для формирования иноязычной лексической компетенции.

Гипотеза исследования: применение мобильного помощника повысит прочность усвоения лексики, если карточки строятся по принципу «одна единица — один концепт» и включают контекст употребления; если упражнения образуют последовательность от узнавания к самостоятельному порождению высказывания; если интервалы повторения рассчитываются индивидуально по результатам ответов; если работа с приложением связана с аудиторными занятиями, а не существует отдельно от них.

Методы исследования: теоретические — анализ методической и психологической литературы, сравнение, обобщение; эмпирические — анкетирование, тестирование, педагогический эксперимент, наблюдение; статистические — обработка результатов, сравнение средних значений, расчёт прироста.

Теоретическую основу работы составили труды по методике обучения иностранным языкам (Н. Д. Гальскова, И. А. Зимняя, Е. И. Пассов, Е. Н. Соловова, А. Н. Щукин), по цифровой дидактике (С. В. Титова, А. П. Авраменко, П. В. Сысоев, Е. С. Полат), а также исследования памяти и учебных стратегий (Г. Эббингауз, А. Пайвио, Г. Рёдигер, Дж. Карпике, П. Вожняк).

Практическая значимость. Разработан прототип мобильного помощника, подготовлен лексический набор из 400 единиц по десяти темам, составлены упражнения пяти типов и методические рекомендации для преподавателя. Материалы могут использоваться на занятиях по иностранному языку в колледже, в самостоятельной работе студентов, а также при подготовке к зачёту или демонстрационному экзамену.

Структура работы. Курсовая работа состоит из введения, двух глав, заключения, списка литературы из 31 наименования и трёх приложений. В первой главе рассматриваются теоретические основания мобильного обучения, психология работы с карточками, устройство упражнений и обзор существующих приложений. Во второй главе описан собственный проект: замысел, требования, структура базы карточек, реализация прототипа и результаты его проверки.

\chapter{Теоретические основы использования мобильных технологий в обучении иностранным языкам}

\section{Мобильное обучение как направление современной методики преподавания иностранных языков}

Термин «мобильное обучение» (mobile learning, m-learning) вошёл в методический обиход сравнительно недавно, но обозначаемое им явление успело сменить несколько форм. В работах начала двухтысячных годов под мобильным обучением понимали использование карманных компьютеров и коммуникаторов для доступа к учебным материалам; сегодня речь идёт прежде всего о смартфоне, который есть практически у каждого студента и который почти всегда включён. С. В. Титова и А. П. Авраменко определяют мобильное обучение как организацию учебного процесса с опорой на мобильные устройства, обеспечивающую непрерывный доступ к учебным ресурсам независимо от места и времени и позволяющую выстроить индивидуальную траекторию работы с материалом [13].

Близкое по смыслу определение даёт словарь методических терминов: мобильное обучение рассматривается как обучение с использованием портативных устройств, объединяющее возможности электронных ресурсов и ситуативность обычного общения [1]. Существенно, что в обоих определениях акцент сделан не на технике как таковой, а на изменении режима учения: материал становится доступен в ту минуту, когда он понадобился, а не только в отведённое расписанием время.

От электронного и дистанционного обучения мобильное отличается рядом признаков, которые принято называть дидактическими свойствами. К ним относят: постоянную доступность (устройство всегда с собой); персонализацию (приложение подстраивается под темп и ошибки конкретного пользователя); интерактивность и немедленную обратную связь (правильность ответа проверяется сразу, без участия преподавателя); мультимедийность (текст, звук, изображение, видео в одном экране); контекстность (возможность учиться в той ситуации, где язык может понадобиться); наконец, возможность фиксировать учебные действия и накапливать статистику [13]. Для преподавателя последнее свойство особенно ценно: появляются данные о том, сколько времени студент действительно работал и какие слова он путает.

Однако сводить мобильное обучение к «обычному курсу, сжатому до размера экрана» было бы ошибкой. Ещё в первых обзорах направления отмечалось, что материалы, перенесённые в телефон без переработки, читаются плохо и почти не используются. Требование переработки не является особенностью цифровой среды: содержание и способы работы с ним всегда отбираются под учебную цель, а не под техническую возможность [3]. Отсюда требование порционности: учебная единица должна укладываться в несколько минут и решаться на одном экране. Эта идея оформилась в отдельное понятие — микрообучение (microlearning), где основной единицей считается небольшой завершённый фрагмент: одно слово с примером, одно правило, одно упражнение из пяти пунктов.

Порционность хорошо согласуется с тем, как реально распределяется внимание студента. Между парами остаётся перемена в десять минут, в дороге — четверть часа, перед сном — ещё немного. Мобильное приложение позволяет занять эти промежутки, тогда как учебник требует сесть за стол и выделить час, которого у первокурсника обычно не находится. Именно поэтому в зарубежной литературе закрепилась формула «learning in the gaps» — учение в промежутках.

Второе важное обстоятельство — распространение модели BYOD (bring your own device), при которой студент пользуется личным устройством, а не техникой учебного заведения. Для колледжа это снимает проблему оснащения компьютерного класса, но ставит вопросы совместимости, информационной безопасности и дисциплины. Практика показывает, что запрет телефонов на занятии не решает проблему отвлечения, а лишь вытесняет его: студент всё равно пользуется устройством, только скрытно. Более продуктивна другая линия — легализовать телефон в отведённых рамках, дав ему учебную роль.

Что даёт мобильный формат применительно к иностранному языку? Прежде всего он удобен для работы над лексикой и аудированием. Оба вида деятельности допускают дробление на короткие операции, не требуют собеседника и легко проверяются автоматически. Слово можно показать, озвучить, попросить выбрать перевод, вписать пропущенную букву, собрать из букв, узнать на слух — всё это реализуется на экране смартфона без потери методического смысла. Гораздо хуже обстоит дело с говорением: диалог требует партнёра и реакции на непредсказуемую реплику, а современные средства автоматического распознавания речи пока ненадёжно оценивают произношение и почти не различают коммуникативную уместность.

Отсюда следует принципиальный вывод, который разделяют большинство исследователей: мобильное приложение не заменяет занятие, а обслуживает его периферию — подготовку к паре, повторение после неё, индивидуальную работу над ошибками. В литературе для такого сочетания используется термин «смешанное обучение» (blended learning), и именно в этой логике строятся удачные примеры внедрения мобильных карточек в вузовскую практику [8].

Нельзя не сказать и об ограничениях, которые приходится учитывать при проектировании. Первое — фрагментарность: короткие сессии дают ощущение движения, но не формируют системы, если их содержание не выстроено по программе. Второе — известный эффект «иллюзии компетентности»: студент уверенно узнаёт слово в списке вариантов и не может употребить его в речи, потому что узнавание и порождение опираются на разные механизмы. Третье — эргономика: маленький экран, неудобная клавиатура, блики на солнце, необходимость держать устройство одной рукой. Четвёртое — отвлечения: уведомление мессенджера обрывает учебную сессию чаще, чем хотелось бы, и приложение должно либо уметь работать в режиме «не беспокоить», либо быть достаточно коротким, чтобы успеть до следующего сигнала.

Перечисленные ограничения не отменяют направления, а задают требования к учебному продукту. Если приложение учитывает программу, чередует узнавание с порождением, объясняет ошибку и укладывается в пять минут, оно превращается из развлечения в рабочий инструмент. Каким образом эти требования конкретизируются в методике работы с карточками и упражнениями, рассматривается в следующих параграфах.

\section{Психолого-педагогические основы работы с лексическими карточками}

Карточка — едва ли не самое старое учебное пособие, дожившее до наших дней без существенных изменений. Это прямоугольный листок, на одной стороне которого написано слово или вопрос, а на другой — перевод, толкование или ответ. Простота конструкции обманчива: карточка воспроизводит в миниатюре полный цикл учебного действия — предъявление, попытку припоминания, проверку и оценку результата. Именно поэтому она легко переводится в цифровую форму: экран заменяет листок, касание заменяет переворачивание, а приложение берёт на себя то, что в бумажном варианте делал сам учащийся, — раскладку карточек по группам и определение момента, когда карточку пора показать снова.

Первая известная система такого рода описана немецким публицистом С. Лейтнером в начале 1970-х годов [17]. Она предельно проста: карточки распределяются по нескольким ящикам; правильный ответ переводит карточку в следующий ящик, ошибка возвращает её в первый. Просматриваются только карточки того ящика, который назначен на сегодня. Чем дальше ящик, тем реже карточка попадается на глаза — и это уже интервальное повторение, только выраженное не в днях, а в порядке очереди.

Обоснование у этой практики появилось задолго до неё. Г. Эббингауз, экспериментируя с заучиванием бессмысленных слогов, установил, что забывание идёт по убывающей кривой: быстрее всего материал теряется в первые часы, затем темп замедляется [15]. В его опытах через час сохранялось около половины заученного, через сутки — примерно треть, а дальше кривая выходила на плато. Для методики отсюда следуют два вывода: повторять надо не «когда забыл», а заранее, на подходе к точке потери; и каждое удачное повторение делает кривую более пологой. Механизмы самой памяти — соотношение непосредственного и опосредованного запоминания, роль смысловых связей — подробно разобраны в работах А. Р. Лурии, и их знание помогает понять, почему бессмысленный список слов удерживается хуже, чем несколько фраз, в которых эти слова работают [6].

Причины забывания не сводятся к простому стиранию следа. Заметную роль играет интерференция: сходные единицы мешают друг другу, и студент, только что выучивший пару «borrow — брать взаймы», упорно подставляет в предложение «lend». Интерференция особенно сильна, когда слова заучиваются списком, без контекста и без различения ситуаций употребления. Отсюда практическое правило, известное преподавателям задолго до появления приложений: не стоит давать в один день десяток слов одной тематической группы, если они близки по значению. Психологические исследования памяти объясняют это так: сходные единицы кодируются в близких «адресах», и при извлечении одна вытесняет другую [2].

Второй механизм, на который опирается работа с карточками, — эффект активного припоминания (testing effect). Суть его в том, что попытка вспомнить сама по себе укрепляет след, причём сильнее, чем повторное чтение того же материала. Г. Рёдигер и Дж. Карпике показали это в серии экспериментов: студенты, которые после первого прочтения проходили тест, через неделю воспроизводили материал заметно лучше тех, кто читал текст повторно, хотя последние были уверены в обратном [22]. Для карточек это ключевое обстоятельство. Карточка ценна не тем, что показывает слово, а тем, что заставляет вспомнить его до того, как ответ откроется. Если пользователь просто листает карточки, проговаривая «да, знаю», эффект близок к нулю.

Третий механизм связан с глубиной переработки материала. Согласно модели уровней обработки, прочность запоминания зависит от того, насколько глубоко информация осмыслена: механическое запоминание формы даёт хрупкий след, а установление связей со смыслом, личным опытом и ситуацией — устойчивый. Практический вывод для разработчика приложения очевиден: карточка с одним лишь переводом работает хуже, чем карточка, где рядом со словом стоит короткий пример, коллокация и, желательно, изображение. Изображение подключает второй канал — зрительный образ, — и это уже ссылка на теорию двойного кодирования А. Пайвио, согласно которой вербальная и образная информация хранятся в разных подсистемах и взаимно усиливают друг друга [20]. Р. Майер, обобщая эксперименты по мультимедийному обучению, сформулировал на этот счёт практическое правило: изображение рядом со словом помогает, если оно поясняет смысл, и мешает, если служит украшением [18].

Существует и более дробная модель, описывающая, что именно значит «знать слово». П. Нейшн различает три стороны: форму (звучание, написание, морфемный состав), значение (понятие, круг сочетаний, коннотации) и употребление (грамматические функции, ограничения, частотность) [19]. Для каждой стороны возможно рецептивное и продуктивное владение. Отсюда следует, что одна и та же карточка не может одинаково хорошо обслуживать все задачи: для узнавания на слух нужна озвучка и перевод, для порождения — русская подсказка и требование написать слово, для употребления — пропуск в предложении. Профессионально сделанный набор карточек поэтому включает несколько типов записей на одну лексическую единицу, а не одну универсальную.

На практике это порождает набор правил, которые стоит перечислить, поскольку они легли в основу нашего проекта:

\begin{enumerate}

  \item на одной карточке — одна единица и один концепт; многозначное слово разводится по отдельным карточкам с пометкой значения;

  \item обязателен контекст: словосочетание или короткое предложение, где слово употреблено типично;

  \item указываются часть речи и, для существительных, форма множественного числа или род, если они неочевидны;

  \item транскрипция и озвучка даются сразу, а не по запросу, — иначе студент закрепляет неверное произношение;

  \item пример должен быть посильным: не более одной незнакомой единицы помимо изучаемой;

  \item карточка не перегружается: два-три поля сверх слова и перевода — предел, за которым экран теряет читаемость.

\end{enumerate}

Отдельный вопрос — направление карточки. Привычная схема «иностранное слово — русский перевод» тренирует узнавание, то есть рецептивное владение. Обратная схема «русское слово — иностранное» требует порождения и потому труднее, но именно она готовит к говорению. Опыт преподавателей показывает: студенты, работавшие только в прямом направлении, узнают слово в тексте, но не могут вызвать его из памяти в нужный момент. Разумное решение — предъявлять карточки в обоих направлениях, но с разной частотой и на разных этапах: сначала прямое, затем обратное, а на продвинутом этапе — пропуск в предложении.

Наконец, нельзя обойти мотивационную сторону. И. А. Зимняя, анализируя психологические механизмы обучения иностранному языку, подчёркивала, что интерес к предмету держится не на внешних требованиях, а на переживании успеха и осмысленности действий [4]. Приложение способно дать и то, и другое: видимый прогресс (сколько слов закреплено, сколько дней подряд занимался) и понятную цель (подготовка к зачёту, к поездке, к профессиональной ситуации). Вместе с тем избыточная внешняя стимуляция — очки, уровни, соревнование — работает недолго и может подменить учебную мотивацию игровой. Об этом подробнее речь пойдёт в следующем параграфе.

Если суммировать, психолого-педагогические основания работы с карточками складываются в три положения. Карточка должна требовать припоминания, а не узнавания постфактум. Повторение должно быть распределённым и подстроенным под индивидуальную кривую забывания. Содержание карточки должно давать контекст и связи, а не только пару «слово — перевод». Эти три положения задают требования и к алгоритму повторений, и к конструкции упражнений.

\section{Упражнения и интервальное повторение в структуре учебного мобильного приложения}

В методике обучения иностранным языкам упражнение понимается как специально организованное действие, выполняемое в учебных условиях с целью формирования навыка или развития умения. Е. И. Пассов построил классификацию, которая до сих пор используется в практике: упражнения делятся на языковые (тренировка формы), условно-речевые (действие по образцу в заданной ситуации) и речевые (самостоятельное высказывание) [7]. Такое деление важно для разработчика: если приложение ограничивается первым типом, оно формирует узнавание, но не готовность к общению. Е. Н. Соловова, разбирая требования к упражнениям в школьном курсе, добавляет к этому ещё один критерий — коммуникативную оправданность: задание должно имитировать речевую задачу, а не проверять память ради памяти [10].

В мобильном формате полный цикл упражнений выстроить труднее, чем в аудитории, поэтому его приходится сокращать до последовательности, сохраняющей логику перехода от формы к смыслу. Наиболее распространённая схема включает четыре шага: знакомство (карточка с озвучкой и примером), узнавание (выбор из вариантов), воспроизведение (ввод слова или сборка из букв) и употребление (пропуск в предложении, перевод фразы, составление высказывания). Переход к следующему шагу открывается после устойчивых правильных ответов, а не по календарю.

Теперь о том, как в приложении рассчитывается момент повторения. Первый алгоритм, получивший широкое распространение, — SM-2, предложенный П. Вожняком в конце 1980-х годов [24]. В нём для каждой карточки хранится коэффициент лёгкости EF, который растёт при верных ответах и падает при ошибках, и интервал, вычисляемый как произведение предыдущего интервала на этот коэффициент. Первые повторения назначаются через один и шесть дней, дальнейшие — по формуле I(n) = I(n–1) × EF. Оценка ответа задаётся шкалой от 0 до 5, но на практике её сводят к четырём кнопкам: «снова», «трудно», «хорошо», «легко».

Современные приложения всё чаще используют более сложные модели, учитывающие индивидуальную кривую забывания пользователя, — например, FSRS, где параметры подбираются по истории ответов. Для учебного проекта, рассчитанного на семестр, такой сложности не требуется: практика показывает, что упрощённый вариант SM-2 с настраиваемыми начальными интервалами даёт результаты, близкие к оптимальным, а главное — предсказуем для преподавателя, который должен понимать, почему карточка вернулась именно сегодня.

Принципиальное значение имеет сам факт распределения. В экспериментах Н. Джепедой и коллегами было показано, что распределённое повторение превосходит концентрированное практически при любых условиях, причём оптимальный интервал зависит от того, на какой срок материал должен сохраниться: чем дальше отложен итоговый тест, тем больше должен быть интервал между занятиями [16]. Для семестрового курса разумным оказывается ориентир, при котором первые повторения идут через день и три, затем через неделю, две и месяц.

Из этого следует практическое требование к объёму сессии. Если интервалы растут, ежедневная нагрузка остаётся небольшой, но она должна быть ежедневной. Приложение, которое показывает пятьдесят карточек за раз, перегружает рабочую память и провоцирует отказ; приложение, которое показывает десять, легко вписывается в перемену. Психологическая литература даёт для этого ориентира известное обоснование — ограниченный объём оперативной памяти, традиционно оцениваемый в семь плюс-минус два элемента, — хотя современные оценки для вербального материала скромнее [11].

Система упражнений в учебном приложении обычно включает несколько типов заданий, и каждый из них решает свою задачу. Выбор перевода из четырёх вариантов тренирует быстрое узнавание, но легко «решается» по форме слова, поэтому дистракторы должны быть правдоподобными — из той же темы, той же части речи, с похожим звучанием. Ввод слова с клавиатуры требует точного воспроизведения формы и потому эффективен для орфографии, но утомляет; его стоит применять к словам, где написание расходится с произношением. Задание с пропуском в предложении (cloze) проверяет употребление и сочетаемость. Сопоставление пар — быстрое упражнение на восстановление связей, удобное для повторения целой темы. Наконец, задание «составь предложение из слов» и перевод с русского на иностранный подводят к продуктивному уровню, хотя автоматическая проверка здесь сложнее и требует списка допустимых вариантов.

Отдельного разговора заслуживает обратная связь. В приложении она должна быть немедленной, короткой и объясняющей. Недостаточно показать красный крестик: студенту нужно увидеть правильный ответ, услышать слово и прочитать пример, иначе ошибка закрепится как случайность. Полезно также фиксировать, какие именно варианты он выбирал, — эта информация позволяет повторно предъявить проблемную пару, а не всю тему целиком.

Геймификация в учебных приложениях — тема, вокруг которой сломано немало копий. Механики вроде серии дней (streak), очков и таблиц лидеров действительно повышают регулярность на коротком отрезке. Однако у них есть побочный эффект: пользователь начинает беречь серию, а не знания, и после её обрыва нередко бросает занятия вовсе — именно такую причину называли студенты в нашем опросе. Более устойчивый вариант — недельная цель по числу сессий и прозрачная статистика освоенных слов, где вознаграждением служит видимое продвижение, а не абстрактный счёт.

Наконец, произношение и аудирование. Синтез речи на современных мобильных платформах достиг качества, достаточного для предъявления слова и короткой фразы, и это делает озвучку обязательным элементом карточки. С распознаванием речи дело сложнее: система уверенно расшифровывает беглую речь носителя, но ошибки в произношении учащегося интерпретирует непоследовательно, а оценка «правильно — неправильно» по одному слову часто несправедлива. Поэтому в учебном приложении разумно ограничиться упражнениями на различение на слух и на повторение за образцом, не выставляя строгой оценки за произношение.

Подведём промежуточный итог. Упражнения и алгоритм повторений — не украшение приложения, а его методическое ядро. Если упражнения выстроены по линии «узнавание — воспроизведение — употребление», а интервалы рассчитываются индивидуально, приложение работает как тренажёр. Если же оно сводится к листанию карточек с переводом, эффект ограничивается кратковременным знакомством с материалом, и рассчитывать на перенос в речь не приходится.

\section{Обзор существующих мобильных решений и требования к учебному помощнику}

Прежде чем проектировать собственное приложение, имеет смысл посмотреть, что уже сделано. Рынок учебных приложений для иностранных языков велик, но при внимательном рассмотрении он распадается на несколько типов продуктов, каждый из которых решает свою задачу и имеет свои слабые места.

Anki — старейший и наиболее уважаемый в преподавательской среде инструмент. Это бесплатное приложение с открытым кодом, работающее на алгоритме SM-2 и поддерживающее синхронизацию между устройствами. Сила Anki в гибкости: пользователь сам создаёт карточки любого вида, включая изображения, звук и клоуз-удаления, а также может подключить готовые колоды, созданные другими. Слабости вытекают из той же гибкости. Интерфейс перегружен настройками, первые полчаса работы уходят на освоение, а не на изучение слов; упражнений, кроме показа карточки, нет; качество готовых колод не проверяется, и студент легко скачивает набор с ошибками [25]. Anki хорош для мотивированного взрослого, который готов потратить время на настройку, и плохо подходит первокурснику, которому нужно начать заниматься сегодня.

Quizlet вырос из сервиса для создания карточек и предлагает несколько режимов работы: просмотр, заучивание, подбор пар, тест. Интерфейс дружелюбен, есть мобильное приложение и веб-версия, наборы легко создаются и распространяются по ссылке. Однако алгоритм повторений в бесплатной версии упрощён, часть функций доступна только по подписке, а реклама отвлекает. Главное же — Quizlet тренирует преимущественно узнавание: задания на выбор и подбор пар преобладают, продуктивных заданий мало [30].

Memrise делает ставку на контекст и на живую речь: к словам прилагаются короткие видео с носителями, а запоминание строится на ассоциациях и повторении в разных режимах. Это сильная сторона — студенты отмечают, что слова запоминаются «с голосом». Ограничения: значительная часть содержания платная, наборы привязаны к собственным курсам, а не к программе колледжа, а упражнения на порождение высказывания почти отсутствуют [28].

Duolingo — самый массовый продукт в этой нише, и его вклад в популяризацию идеи ежедневных коротких занятий трудно переоценить. Геймификация здесь доведена до совершенства: серии, жизни, лиги, уведомления. Но Duolingo предлагает линейный курс, а не инструмент для работы с конкретным списком слов; объяснения грамматики минимальны, а лексика подбирается под сценарий курса, а не под профессиональную программу СПО. Кроме того, реклама и платная подписка встроены в сам процесс обучения [26].

Русскоязычные сервисы — Lingualeo и Puzzle English — выросли из работы с текстами и аудированием. Lingualeo предлагает чтение и просмотр видео с подсказками перевода и добавлением слов в личный словарь, Puzzle English специализируется на тренировке восприятия речи на слух. Оба сервиса полезны как источник материала, но работа с лексикой в них вторична по отношению к контенту: словарь пополняется по ходу чтения, система повторений упрощена, а методической привязки к учебному плану нет [27; 29].

Сравнение по формальным признакам приведено в таблице 1. Оценки выставлены по результатам самостоятельного тестирования бесплатных версий приложений в сентябре 2025 года и отражают точку зрения учебного применения, а не качество продукта в целом.

\begin{table}[htbp]
\centering
\footnotesize
\caption{Сравнение мобильных приложений для работы с иноязычной лексикой}
\begin{tabular}{|p{3.0cm}|p{1.9cm}|p{2.2cm}|p{1.9cm}|p{1.5cm}|p{1.5cm}|p{2.2cm}|}
\hline
Приложение & Свои карточки & Интервальное повторение & Упражнения & Озвучка & Офлайн & Русский интерфейс \\
\hline
Anki & да & да (SM-2) & нет & да & да & частично \\
\hline
Quizlet & да & частично & да* & да & нет & да \\
\hline
Memrise & нет & да & частично & да & частично & да \\
\hline
Duolingo & нет & внутр. & да & да & частично & да \\
\hline
Lingualeo & да & частично & частично & да & нет & да \\
\hline
Puzzle English & да & нет & да* & да & частично & да \\
\hline
\end{tabular}
\end{table}

* Задания Quizlet и Puzzle English построены преимущественно на узнавании; «внутр.» — собственный алгоритм курса, не настраиваемый пользователем.

Из таблицы видно, что ни один из продуктов не закрывает все позиции одновременно. Anki силён в алгоритме, но беден на упражнения; Quizlet удобен, но слаб в повторениях; массовые курсы не привязаны к программе. Складывается ниша для небольшого приложения, которое соединило бы три вещи: карточки, построенные по программе конкретной дисциплины; упражнения, доводящие до порождения; и прозрачный алгоритм повторений, понятный и студенту, и преподавателю.

На основании анализа литературы и обзора продуктов были сформулированы требования к учебному помощнику. Функциональные требования:

\begin{enumerate}

  \item работа с собственными наборами карточек, привязанными к темам рабочей программы;

  \item карточка содержит слово, транскрипцию, перевод, часть речи, пример употребления и озвучку;

  \item расчёт даты следующего повторения по результатам ответа (упрощённый SM-2);

  \item не менее пяти типов упражнений, включая продуктивные;

  \item немедленная обратная связь с показом правильного ответа и примера;

  \item статистика: освоенные слова, точность ответов, серии занятий;

  \item напоминания в заданное время, отключаемые одной кнопкой;

  \item экспорт и импорт набора в открытом формате для обмена между группами.

\end{enumerate}

Эргономические требования вытекали из условий использования: одна рука, экран 5–6 дюймов, яркое освещение или, наоборот, темнота в транспорте. Отсюда — крупные кнопки, минимальный ввод с клавиатуры, читаемый шрифт, тёмная тема, отсутствие мелких элементов в нижней части экрана, где лежит большой палец. Технические требования: работа без подключения к сети, размер дистрибутива до 30 МБ, поддержка Android 8.0 и выше, отсутствие сбора персональных данных, поскольку приложение используется несовершеннолетними.

Перечисленные требования и определили облик проекта, описанного во второй главе. Стоит подчеркнуть, что они не выводятся один из другого автоматически: между бытовым запросом «удобно в транспорте» и методическим требованием «доводить до порождения» лежит выбор, который делает разработчик. Наш выбор состоял в том, чтобы пожертвовать универсальностью ради предсказуемости: небольшой набор, жёсткая последовательность упражнений, простой и объяснимый алгоритм повторений. Проигрыш в гибкости компенсируется тем, что преподаватель понимает, что именно происходит с его студентом, а студент за пять минут видит результат своей работы.

Отдельно стоит объяснить, почему выбрана собственная разработка, а не адаптация готового продукта. Готовое приложение дешевле и быстрее, но требует от преподавателя отказаться от части методических решений: нельзя изменить алгоритм повторений, трудно ограничить набор темой занятия, почти невозможно получить данные о работе группы. Собственная разработка, даже в виде прототипа, возвращает эти решения преподавателю — в этом её главный смысл. Цена — время на проектирование, программирование и сопровождение, а также зависимость от одного разработчика. Для кафедры, где учатся будущие программисты, такой обмен выглядит оправданным: приложение становится не только средством обучения языку, но и учебным проектом.

\chapter{Разработка мобильного помощника и его проверка}

\section{Замысел, целевая аудитория и функциональные требования к приложению}

Разработанный помощник получил рабочее название Lexi. Название короткое, легко читается на иконке и не обещает больше, чем приложение умеет: это именно помощник, а не курс и не репетитор. Замысел состоял в том, чтобы создать компактный инструмент, который закрывает ежедневную рутину работы над словом и при этом остаётся частью учебного процесса, а не отдельным развлечением.

Цель приложения сформулирована так: обеспечить студенту возможность ежедневно, в течение пяти–семи минут, повторять лексику текущей темы в режиме активного припоминания, а преподавателю — получать обратную связь о том, какие слова группа усвоила и какие требуют возвращения на занятии.

Целевая аудитория — студенты первого и второго курсов колледжа, изучающие английский язык на уровне A2–B1. Возраст 16–19 лет, цифровая грамотность высокая, бюджет на подписки нулевой, устройство преимущественно на Android (по нашему опросу — 43 из 46 человек). Второй пользователь, косвенный, — преподаватель: он формирует набор, видит сводную статистику и решает, что вынести на аудиторную работу.

Чтобы не проектировать «для всех сразу», были составлены два сценария. Первый: Анна, 17 лет, первый курс, готовится к лексическому диктанту в пятницу. Ей нужно за три дня повторить около восьмидесяти слов темы «Повседневная жизнь», причём она не помнит, какие именно слова забыла. Приложение показывает очередь на сегодня, состоящую из двадцати карточек, и по итогам сессии формирует список проблемных единиц, который Анна просматривает перед диктантом. Второй сценарий: Дмитрий, 18 лет, едет в автобусе двадцать минут. Он открывает приложение без конкретной цели, проходит сессию из десяти карточек с упражнениями, получает напоминание на завтра и закрывает его.

Из сценариев выросли пользовательские задачи: быстро увидеть, что делать сегодня; потратить на сессию не больше десяти минут; понять, почему ответ неверен; увидеть, что прогресс есть; не потерять данные при смене устройства. Формулировки этих задач легли в основу функциональных требований, которым присвоены идентификаторы, чтобы в дальнейшем можно было проверить их выполнение. При проектировании учитывались и общие дидактические требования к учебным пособиям — наглядность, доступность, систематичность и сознательность усвоения: приложение не должно превращаться в «чёрный ящик», выдающий оценки без объяснения [5].

Функциональные требования к прототипу:

\begin{enumerate}

  \item FR-1. Приложение запускается без авторизации и без подключения к сети; при первом запуске загружается встроенный набор из десяти тем.

  \item FR-2. Главный экран показывает план на день: число карточек к повторению, число новых слов, ориентировочное время сессии.

  \item FR-3. Карточка отображает слово, транскрипцию, часть речи, перевод, пример и кнопку озвучки; ответ открывается касанием, после чего пользователь выбирает одну из четырёх оценок.

  \item FR-4. После знакомства с карточкой предлагается упражнение одного из пяти типов; тип выбирается автоматически с учётом истории ответов.

  \item FR-5. Оценка ответа пересчитывает интервал следующего повторения и коэффициент лёгкости карточки.

  \item FR-6. Экран статистики показывает число освоенных слов по темам, точность за неделю и календарь занятий.

  \item FR-7. Напоминание приходит в выбранное пользователем время, если за день не было ни одной сессии; напоминание отключается в настройках.

  \item FR-8. Набор можно выгрузить в файл CSV и загрузить обратно, в том числе набор другой группы.

  \item FR-9. Тёмная тема включается вручную или по системной настройке.

\end{enumerate}

Нефункциональные требования касались в основном ограничений среды. Приложение должно работать на устройствах с Android 8.0 и выше без сетевого подключения, занимать не более 30 МБ вместе со звуковыми файлами, запускаться быстрее чем за две секунды на бюджетном устройстве и не содержать аналитических библиотек, передающих данные третьим лицам. Последнее требование не техническое, а этическое: пользователи — несовершеннолетние, и собирать их данные без явной необходимости не следует.

Архитектура прототипа трёхслойная и намеренно простая. Слой представления построен на Jetpack Compose — декларативном инструменте интерфейса, который позволяет описать экран как функцию от состояния [21]. Слой данных опирается на локальную базу Room, где хранятся карточки, темы, история ответов и параметры повторений. Слой логики содержит планировщик повторений, генератор упражнений и обработчик статистики. Учебный контент поставляется вместе с приложением в виде JSON-файла, что позволяет обновлять набор без пересборки программы.

Границы прототипа были заданы сразу, чтобы не расползаться по объёму. В первую версию не вошли: версия для iOS, синхронизация между устройствами, социальные функции, автоматическая генерация карточек из текста и распознавание речи. Эти возможности перечислены в перспективах развития, но их реализация требует либо серверной части, либо существенно большего времени, чем отведено на курсовую работу.

Таким образом, замысел свёлся к небольшому, но методически последовательному продукту: карточка, упражнение, оценка, пересчёт интервала — и всё это в пределах пяти минут и без интернета. Конкретное наполнение этого каркаса — структура базы карточек и система упражнений — описано в следующем параграфе.

\section{Структура базы карточек и система упражнений}

Содержательное ядро приложения — база карточек. От того, как она устроена, зависит не меньше, чем от алгоритма повторений: неудачно составленная карточка портит впечатление от самой продуманной программы. Поэтому проектирование началось не с экранов, а с модели данных и правил заполнения.

Модель включает четыре сущности. Карточка (Card) описывает одну лексическую единицу и хранит поля, перечисленные в таблице 2. Тема (Deck) объединяет карточки по программе и содержит название, уровень и порядковый номер. Запись о повторении (Review) фиксирует дату, оценку, тип упражнения и время ответа. Попытка (Attempt) хранит выбранный вариант и правильный ответ — это нужно для анализа типичных ошибок, а не только для статистики.

\begin{table}[htbp]
\centering
\small
\caption{Поля карточки лексической единицы}
\begin{tabular}{|p{3.2cm}|p{3.0cm}|p{9.6cm}|}
\hline
Поле & Тип & Назначение \\
\hline
id & целое & уникальный идентификатор \\
\hline
word & строка & слово или устойчивое сочетание \\
\hline
transcription & строка & транскрипция в формате МФА \\
\hline
part\_of\_speech & перечисление & существительное, глагол, прилагательное, наречие, фраза \\
\hline
translation & строка & основное значение на русском языке \\
\hline
example & строка & пример употребления с переводом \\
\hline
collocation & строка & типичное словосочетание \\
\hline
audio & файл & озвучка слова и примера \\
\hline
image & файл & изображение, если оно помогает различению \\
\hline
deck\_id & целое & принадлежность к теме \\
\hline
ef & вещественное & коэффициент лёгкости (по умолчанию 2,5) \\
\hline
interval & целое & текущий интервал в днях \\
\hline
due & дата & дата следующего повторения \\
\hline
reps & целое & число успешных повторений подряд \\
\hline
\end{tabular}
\end{table}

Лексический набор охватывает десять тем, соответствующих рабочей программе дисциплины «Иностранный язык» для первого курса: повседневная жизнь, дом и быт, еда и покупки, здоровье, транспорт и путешествия, учёба и профессия, деловая коммуникация, технологии, экология, свободное время. По каждой теме подготовлено по сорок единиц, всего четыреста. Распределение по частям речи получилось неравномерным и отражает реальную структуру частотного списка: существительные — около 55 \%, глаголы — 25 \%, прилагательные — 12 \%, наречия и устойчивые сочетания — 8 \%.

Отбор слов проводился по трём критериям. Первый — частотность: приоритет отдавался единицам, входящим в первые две тысячи наиболее употребительных слов английского языка, поскольку именно они обеспечивают понимание основного массива текстов [23]. Второй — соответствие программе: набор сверялся с лексическим минимумом дисциплины, чтобы приложение не расходилось с тем, что преподаватель даёт на занятии. Третий — профессиональная направленность: для тем «Деловая коммуникация» и «Технологии» подбирались единицы, встречающиеся в типовых ситуациях будущей специальности (собеседование, деловая переписка, инструкция, презентация).

Отдельная забота — так называемые ложные друзья переводчика и интернационализмы. Слова вроде magazine («журнал», а не «магазин»), accurate («точный», а не «аккуратный»), data («данные», а не «дата») вынесены в отдельные карточки с предупреждением в поле примера, потому что именно они дают устойчивые ошибки, которые не исправляются простым повторением.

Формулировки внутри карточки подчинены правилу «одна единица — один концепт». Если слово многозначно, оно разводится на две карточки с пометкой значения: bill — «счёт» и bill — «законопроект». Такой подход увеличивает число карточек, но снижает путаницу: студент не пытается удержать в одной записи два несвязанных смысла. Пример употребления всегда короче двенадцати слов и содержит только одну незнакомую единицу — саму изучаемую.

Вот как выглядит карточка темы «Деловая коммуникация» в готовом виде: слово deadline, транскрипция {\ipa [ˈdedlaɪn]}, часть речи — существительное, перевод — «крайний срок», пример — “The deadline for the report is Friday”, коллокации — meet a deadline, miss a deadline, озвучка — два файла, для слова и для примера. Такая карточка даёт форму, значение и употребление одновременно, а поле коллокаций подсказывает, с чем слово сочетается.

Система упражнений построена как лестница из пяти ступеней, и каждая ступень опирается на предыдущую. Первая — выбор перевода из четырёх вариантов: самое лёгкое задание, которое вводит слово в работу и даёт быстрый успех. Вторая — выбор иностранного слова по русскому переводу, то есть обратное направление, уже требующее припоминания формы. Третья — ввод слова с клавиатуры с подсказкой первой буквы: тренирует орфографию, но не наказывает за одну опечатку, если слово узнано верно. Четвёртая — пропуск в предложении, где нужно выбрать или вписать слово по смыслу: здесь проверяется сочетаемость. Пятая — сборка предложения из перемешанных слов и перевод короткой фразы с русского: продуктивный уровень, на котором приложение проверяет не точное совпадение, а наличие ключевых единиц в ответе.

Технически упражнения генерируются из той же карточки, а не пишутся отдельно. Генератор берёт слово, пример и коллокации и строит задание выбранного типа: для cloze вырезает слово из примера, для выбора вариантов подбирает три дистрактора, для сборки предложения перемешивает слова примера. Дистракторы подбираются не случайно, а из той же темы и той же части речи, с учётом длины и, по возможности, сходства звучания — иначе правильный ответ угадывается по форме.

Переход между ступенями не привязан к календарю. Карточка считается освоенной на текущем уровне, если пользователь дал три верных ответа подряд, причём хотя бы один — в упражнении другого типа. После этого она переходит на следующую ступень, а при двух ошибках подряд возвращается на ступень ниже. Такая схема, заимствованная из практики адаптивных тренажёров, не даёт «провалиться» в бесконечное повторение одного и того же и в то же время не позволяет пройти тему на узнавании.

Ошибки не наказываются, а собираются. Если пользователь дважды ошибся в одном и том же слове, карточка попадает в отдельный список «трудные слова», который можно открыть перед зачётом. Преподавателю этот список полезен не меньше: он показывает, какие единицы стоит разобрать на занятии, — и это как раз тот случай, когда приложение экономит аудиторное время, а не отнимает его.

\section{Реализация прототипа и сценарий учебной работы}

Прототип реализован для платформы Android. Выбор платформы определился результатами опроса: подавляющее большинство студентов пользуется устройствами на Android, а разработка под iOS потребовала бы другого инструментария и, главное, другого бюджета. Средой разработки стал Android Studio, языком — Kotlin, интерфейс построен на Jetpack Compose, хранение данных — на Room поверх SQLite. Минимальная поддерживаемая версия системы — Android 8.0, целевая — Android 14.

Приложение состоит из шести экранов. Главный показывает план на день: сколько карточек ждёт повторения, сколько новых слов предлагается и сколько времени займёт сессия. Экран темы открывает список колод с прогрессом по каждой. Экран карточки — основной рабочий экран: слово, транскрипция, часть речи, кнопка озвучки, а после касания — перевод, пример и четыре кнопки оценки. Экран упражнения отличается в зависимости от типа задания. Экран статистики показывает число освоенных слов по темам, точность ответов за неделю и календарь занятий. Экран настроек содержит время напоминания, размер дневной порции, тему оформления и импорт-экспорт набора.

Планировщик повторений реализует упрощённый SM-2. Для каждой карточки хранятся коэффициент лёгкости EF (по умолчанию 2,5), текущий интервал I и число успешных повторений подряд n. При оценке «снова» интервал сбрасывается на один день, а EF уменьшается на 0,2, но не опускается ниже 1,3. При оценке «трудно» интервал умножается на 1,2. При оценке «хорошо» — на EF. При оценке «легко» — на EF с добавлением 0,15 к самому коэффициенту. Первые два повторения назначаются через один и шесть дней, как в оригинальном алгоритме, дальнейшие рассчитываются по формуле.

Озвучка построена на системном синтезаторе речи Android. Это решение имеет свои издержки: качество голоса зависит от устройства и установленных языковых пакетов, а на некоторых прошивках британский акцент подменяется американским. Тем не менее отказ от записи живых дикторов позволил сохранить размер приложения в пределах 18 МБ и обойтись без серверной части. Звук кэшируется после первого воспроизведения, поэтому повторные сессии не требуют сети и не расходуют заряд.

Напоминания реализованы через WorkManager — планировщик отложенных задач, который учитывает режим энергосбережения и не срабатывает, если пользователь уже занимался сегодня [31]. Важная деталь: уведомление не должно быть навязчивым. В первой версии оно приходило в фиксированное время и раздражало; после доработки оно появляется только при отсутствии сессии за день и формулируется нейтрально («Осталось 12 карточек, около четырёх минут»), без упрёков и восклицательных знаков.

Проверка прототипа шла в два этапа. Сначала — юзабилити-тестирование с двенадцатью студентами, каждый прошёл три сессии с комментариями вслух. Выявились четыре проблемы: кнопка «снова» стояла рядом с «хорошо» и нажималась случайно; смысл оценки «трудно» был неочевиден; озвучка молчала в беззвучном режиме; после сессии было непонятно, что делать дальше. Кнопки разнесли, к «трудным» добавили подсказку и предупреждение о звуке, а на экран завершения вынесли кнопку «повторить трудные».

Второй этап — педагогическая проверка, о которой речь пойдёт в следующем параграфе. Здесь же стоит описать сценарий учебной работы, поскольку без него приложение остаётся набором карточек, а не методикой. Сценарий рассчитан на семестр и включает три элемента.

Первый — пятиминутная входная сессия в начале занятия. Студенты открывают приложение и проходят очередь повторения по прошлой теме, пока преподаватель проверяет присутствие и готовит материалы. Преподаватель видит на своём экране сводку: какие слова группа путает чаще всего. Эти слова выносятся на доску и разбираются в течение двух-трёх минут.

Второй элемент — работа с новой лексикой на занятии, где приложение играет вспомогательную роль: студенты находят в нём карточки новой темы, слушают озвучку, проверяют транскрипцию. Здесь важно, чтобы приложение не подменяло объяснение: значение слова всё равно раскрывается в контексте, в диалоге, в упражнении с преподавателем.

Третий элемент — домашнее задание. Оно формулируется не как «выучите слова», а как «пройдите сессию в приложении три раза до пятницы и разберите трудные слова». Такая формулировка измерима, и её выполнение видно по статистике. Домашнее задание занимает десять минут, а не час, и это принципиально: студент, которому задают час, не делает ничего, студент, которому задают десять минут, делает их в транспорте.

Из опыта сопровождения групп сложились методические рекомендации преподавателю:

\begin{enumerate}

  \item вводить приложение не с первого занятия, а когда группа освоилась;

  \item начинать с одной темы и двадцати карточек, а не с полного набора — объём отпугивает;

  \item показывать на занятии, как выглядит сессия, один раз, всем вместе, на большом экране;

  \item не проверять статистику как дневник: она нужна для помощи, а не контроля;

  \item раз в две недели выносить «трудные слова» группы на отдельное упражнение в аудитории;

  \item оставлять студенту выбор: приложение дополняет занятие, но оценка не должна зависеть только от него.

\end{enumerate}

Ограничения прототипа перечислены, чтобы не создавать ложного впечатления о готовом продукте: нет версии для iOS и синхронизации между устройствами, набор ограничен четырьмястами единицами, а проверка продуктивных ответов сводится к наличию ключевых единиц. Это направления дальнейшей работы, а не скрытые дефекты.

\section{Оценка эффективности и перспективы развития помощника}

Педагогическая проверка проводилась в сентябре — ноябре 2025 года на базе одной из образовательных организаций среднего профессионального образования. В ней участвовали 24 студента первого курса, разделённые на две группы по двенадцать человек. Разделение было не случайным в строгом смысле: группы формировались по учебным подгруппам, уже существующим в расписании, — и это ограничение придётся учитывать при интерпретации. Общая логика проверки соответствовала принятой в педагогике схеме эксперимента: констатирующий этап, формирующий этап и контрольный срез с последующей обработкой результатов [9].

Обе группы занимались по одной программе, с одним преподавателем и в одинаковом объёме часов. Экспериментальная группа дополнительно работала с приложением: пять минут в начале занятия и домашние сессии три раза в неделю. Контрольная группа получала те же слова в бумажном виде — списком с переводом и примерами — и выполняла те же домашние задания, но без приложения.

До начала работы был проведён входной срез: 50 слов на узнавание (выбор перевода) и 20 слов на продукцию (перевод с русского на английский). Слова отбирались из тем, которые ещё не изучались, поэтому низкие результаты ожидались и подтвердились: обе группы узнавали около трети единиц, а продуцировали — около одной седьмой. Разница между группами на входе была незначительной.

Итоговый срез проводился через неделю после завершения темы и включал те же 50 и 20 единиц, но в другом порядке и с другими дистракторами. Кроме того, через две недели после итогового среза был проведён отсроченный тест на узнавание — он должен был показать, сохранился ли материал, когда эффект подготовки к проверке перестал действовать. Результаты приведены в таблице 3.

\begin{table}[htbp]
\centering
\small
\caption{Результаты входного, итогового и отсроченного срезов, \% выполнения}
\begin{tabular}{|p{4.4cm}|p{2.2cm}|p{2.2cm}|p{2.2cm}|p{2.2cm}|p{2.6cm}|}
\hline
Показатель & КГ, вход & КГ, итог & ЭГ, вход & ЭГ, итог & ЭГ, отсроченный \\
\hline
Узнавание (50 единиц) & 34 & 61 & 36 & 78 & 71 \\
\hline
Продукция (20 единиц) & 12 & 34 & 14 & 52 & — \\
\hline
Точность в упражнениях & — & 68 & — & 84 & — \\
\hline
\end{tabular}
\end{table}

Прирост в экспериментальной группе оказался заметно выше. По узнаванию разница между группами составила 17 процентных пунктов, по продукции — 18. Ещё показательнее отсроченный срез: через две недели экспериментальная группа сохранила 71 \% против 52 \% в контрольной, то есть преимущество не только не исчезло, но и увеличилось в относительном выражении. Это согласуется с тем, что известно об интервальном повторении: оно работает не столько на момент проверки, сколько на длительную перспективу.

Точность в упражнениях внутри приложения (84 \%) выше, чем в контрольной группе при работе с бумажными заданиями (68 \%), и это ожидаемо: приложение проверяет сразу, а бумажное задание проверяется на следующем занятии, когда ошибка уже закрепилась.

Дополнительно было проведено анкетирование участников экспериментальной группы. Двадцать два человека из двадцати четырёх оценили работу с приложением как удобную, восемнадцать отметили, что стали заниматься регулярнее, шестеро пожаловались на качество синтезированной речи, четверо просили добавить больше продуктивных упражнений. Последнее пожелание показательно: студенты воспринимают приложение не как учебник, а как тренажёр, где уместна игровая составляющая, но лишь до определённой границы.

Полученные результаты нельзя переоценивать. Выборка мала, группы не были уравнены случайным образом, преподаватель знал, какая из них экспериментальная, а сам эксперимент длился шесть недель — слишком мало для устойчивых выводов. Возможен и эффект новизны: часть прироста объясняется повышенным вниманием к предмету. Тем не менее направление эффекта устойчиво, а его величина достаточно велика, чтобы оправдать дальнейшую работу.

Перспективы развития видятся в нескольких направлениях: версия для iOS и синхронизация прогресса между устройствами; переход от упрощённого SM-2 к адаптивной модели, подбирающей параметры по истории ответов; генерация карточек из учебного текста, когда приложение само выделяет незнакомые единицы; распознавание речи для тренировки произношения без строгой оценки; наконец, панель преподавателя со сводной статистикой по группе, чтобы «трудные слова» выявлялись автоматически.

Отдельно стоит сказать о рисках, которые проявились в ходе работы. Приложение экономит время, но может создать у студента ощущение, что занятий в нём достаточно. Опыт шести недель показал: те, кто ограничивался только сессиями, прибавили меньше остальных. Это подтверждает тезис первой главы: мобильный помощник обслуживает периферию учебного процесса, а его ядро остаётся за преподавателем и живым общением.

\chapter*{Заключение}
\addcontentsline{toc}{chapter}{Заключение}
\markboth{Заключение}{Заключение}

Проведённая работа началась с практической трудности, знакомой каждому преподавателю языка: слова не удерживаются в памяти, если с ними не работать между занятиями, а времени на эту работу в сетке часов нет. Мобильный помощник, разработанный и проверенный в ходе исследования, — попытка закрыть именно эту щель, не подменяя собой ни занятие, ни учебник.

Первая задача состояла в изучении сущности мобильного обучения. Анализ литературы показал, что его нельзя свести к переносу учебных материалов на экран телефона. Мобильное обучение меняет режим работы с материалом: делает его доступным в момент необходимости, дробит на короткие порции, обеспечивает немедленную обратную связь и накапливает данные о действиях учащегося. Одновременно выявились и ограничения — фрагментарность, риск иллюзии компетентности, эргономические неудобства, отвлечения. Отсюда принципиальный вывод: приложение обслуживает периферию учебного процесса, а его ядро остаётся в аудитории.

Вторая задача касалась психолого-педагогических оснований работы с карточками. Удалось показать, что карточка эффективна не сама по себе, а благодаря трём механизмам: активному припоминанию, распределённому повторению и глубине переработки материала. Из этого выросли конкретные правила составления карточек — одна единица на карточку, обязательный контекст, транскрипция и озвучка, ограниченное число полей. Эти правила легли в основу лексического набора, а не остались теоретической декларацией.

Третья задача была связана с упражнениями и алгоритмами повторения. Рассмотренные классификации упражнений и модели интервального повторения (от системы Лейтнера до SM-2 и современных адаптивных алгоритмов) позволили выстроить последовательность из пяти ступеней — от выбора перевода до перевода фразы — и реализовать упрощённый SM-2, понятный и студенту, и преподавателю. Такой выбор может показаться консервативным на фоне более сложных моделей, но для семестрового курса предсказуемость важнее математической изощрённости.

Четвёртая задача — сравнение существующих приложений — показала, что рынок закрывает отдельные потребности, но не их комбинацию. Anki силён в алгоритме и слаб в упражнениях, Quizlet удобен, но тренирует преимущественно узнавание, массовые курсы не привязаны к программе колледжа. Это и определило нишу собственного проекта и перечень требований к нему.

Пятая задача носила практический характер. Был разработан прототип приложения Lexi для Android: база из 400 карточек по десяти темам, пять типов упражнений, планировщик повторений, статистика и напоминания. Прототип прошёл юзабилити-тестирование, в ходе которого были исправлены четыре существенные ошибки интерфейса, а затем педагогическую проверку в течение шести недель.

Результаты проверки подтверждают гипотезу, хотя и с оговорками. Прирост по узнаванию в экспериментальной группе составил 42 процентных пункта против 27 в контрольной, по продукции — 38 против 22. Через две недели после итогового среза экспериментальная группа сохранила 71 \% узнавания, контрольная — 52 \%. Разница устойчива и по направлению, и по величине, однако малая выборка, отсутствие случайного распределения и эффект новизны не позволяют считать вывод окончательным. Корректнее говорить о подтверждении гипотезы в первом приближении и о необходимости более длительной проверки.

Практическая значимость работы состоит в трёх продуктах, которые могут использоваться непосредственно: прототип приложения, лексический набор из 400 единиц и методические рекомендации преподавателю. Кроме того, описанный сценарий учебной работы — пятиминутная сессия в начале занятия, работа с новыми словами в аудитории, десятиминутное домашнее задание — переносим на другие дисциплины, где требуется запоминание большого объёма единиц: терминологии специальности, профессиональной лексики, нормативных формулировок.

Дальнейшая работа видится в нескольких направлениях: версия для iOS и синхронизация прогресса, переход к адаптивной модели повторений, генерация карточек из учебных текстов, тренировка произношения с распознаванием речи, панель преподавателя со сводной статистикой по группе. Каждое из них требует отдельного проекта, и ни одно не отменяет главного вывода.

Этот вывод прост и, возможно, не нов: технология не учит вместо преподавателя. Приложение может напомнить, показать, проверить и посчитать, но мотив учиться, понимание смысла усилий и умение пользоваться языком в реальной ситуации формируются в общении. Мобильный помощник полезен ровно настолько, насколько он встроен в это общение — и именно в таком качестве он рассматривается в настоящей работе.

\chapter*{Список литературы}
\addcontentsline{toc}{chapter}{Список литературы}
\markboth{Список литературы}{Список литературы}

\begin{thebibliography}{99}

\bibitem{} Азимов, Э. Г. Новый словарь методических терминов и понятий (теория и практика обучения языкам) / Э. Г. Азимов, А. Н. Щукин. — Москва : ИКАР, 2009. — 448 с. — Текст : непосредственный.

\bibitem{} Андерсон, Дж. Когнитивная психология / Дж. Андерсон. — 5-е изд. — Санкт-Петербург : Питер, 2002. — 496 с. — Текст : непосредственный.

\bibitem{} Гальскова, Н. Д. Теория обучения иностранным языкам. Лингводидактика и методика : учебное пособие / Н. Д. Гальскова, Н. И. Гез. — 8-е изд., стер. — Москва : Академия, 2017. — 336 с. — Текст : непосредственный.

\bibitem{} Зимняя, И. А. Психология обучения иностранным языкам в школе / И. А. Зимняя. — Москва : Просвещение, 1991. — 222 с. — Текст : непосредственный.

\bibitem{} Коджаспирова, Г. М. Педагогика в схемах и таблицах : учебное пособие / Г. М. Коджаспирова. — Москва : Проспект, 2016. — 248 с. — Текст : непосредственный.

\bibitem{} Лурия, А. Р. Лекции по общей психологии / А. Р. Лурия. — Санкт-Петербург : Питер, 2006. — 320 с. — Текст : непосредственный.

\bibitem{} Пассов, Е. И. Основы коммуникативной методики обучения иноязычному общению / Е. И. Пассов. — Москва : Русский язык, 1989. — 276 с. — Текст : непосредственный.

\bibitem{} Полат, Е. С. Новые педагогические и информационные технологии в системе образования : учебное пособие / Е. С. Полат, М. Ю. Бухаркина. — Москва : Академия, 2008. — 272 с. — Текст : непосредственный.

\bibitem{} Сластенин, В. А. Педагогика : учебник для студентов педагогических учебных заведений / В. А. Сластенин, И. Ф. Исаев, Е. Н. Шиянов. — Москва : Академия, 2013. — 576 с. — Текст : непосредственный.

\bibitem{} Соловова, Е. Н. Методика обучения иностранным языкам : базовый курс лекций / Е. Н. Соловова. — Москва : Просвещение, 2006. — 239 с. — Текст : непосредственный.

\bibitem{} Солсо, Р. Когнитивная психология / Р. Солсо. — 6-е изд. — Санкт-Петербург : Питер, 2006. — 589 с. — Текст : непосредственный.

\bibitem{} Сысоев, П. В. Информационные и коммуникационные технологии в лингвистическом образовании : учебное пособие / П. В. Сысоев. — Москва : Либроком, 2013. — 264 с. — Текст : непосредственный.

\bibitem{} Титова, С. В. Мобильное обучение иностранным языкам : учебное пособие / С. В. Титова, А. П. Авраменко. — Москва : ИКАР, 2014. — 224 с. — Текст : непосредственный.

\bibitem{} Щукин, А. Н. Обучение иностранным языкам: теория и практика : учебное пособие / А. Н. Щукин. — Москва : Филоматис, 2006. — 480 с. — Текст : непосредственный.

\bibitem{} Эббингауз, Г. Основы психологии / Г. Эббингауз ; перевод с немецкого. — Санкт-Петербург : Общественная польза, 1912. — 304 с. — Текст : непосредственный.

\bibitem{} Cepeda, N. J. Distributed practice in verbal recall tasks: a review and quantitative synthesis / N. J. Cepeda, H. Pashler, E. Vul, J. T. Wixted, D. Rohrer // Psychological Bulletin. — 2006. — Vol. 132, № 3. — P. 354–380. — Текст : непосредственный.

\bibitem{} Leitner, S. So lernt man lernen: der Weg zum Erfolg / S. Leitner. — Freiburg : Herder, 1972. — 224 S. — Текст : непосредственный.

\bibitem{} Mayer, R. E. Multimedia learning / R. E. Mayer. — 2nd ed. — Cambridge : Cambridge University Press, 2009. — 304 p. — Текст : непосредственный.

\bibitem{} Nation, I. S. P. Learning vocabulary in another language / I. S. P. Nation. — Cambridge : Cambridge University Press, 2001. — 496 p. — Текст : непосредственный.

\bibitem{} Paivio, A. Mental representations: a dual coding approach / A. Paivio. — New York : Oxford University Press, 1986. — 322 p. — Текст : непосредственный.

\bibitem{} Phillips, B. Android. Программирование для профессионалов / B. Phillips, C. Stewart, K. Marsicano. — Санкт-Петербург : Питер, 2019. — 704 с. — Текст : непосредственный.

\bibitem{} Roediger, H. L. Test-enhanced learning: taking memory tests improves long-term retention / H. L. Roediger, J. D. Karpicke // Psychological Science. — 2006. — Vol. 17, № 3. — P. 249–255. — Текст : непосредственный.

\bibitem{} Schmitt, N. Vocabulary in language teaching / N. Schmitt. — Cambridge : Cambridge University Press, 2000. — 224 p. — Текст : непосредственный.

\bibitem{} Wozniak, P. A. Optimization of learning : master's thesis / P. A. Wozniak. — Poznań : University of Technology, 1990. — 154 p. — Текст : непосредственный.

\bibitem{} Anki: официальный сайт программы для интервального повторения. — URL: https://apps.ankiweb.net (дата обращения: 14.09.2025). — Текст : электронный.

\bibitem{} Duolingo: официальный сайт образовательной платформы. — URL: https://www.duolingo.com (дата обращения: 14.09.2025). — Текст : электронный.

\bibitem{} Lingualeo: официальный сайт сервиса для изучения английского языка. — URL: https://lingualeo.com (дата обращения: 16.09.2025). — Текст : электронный.

\bibitem{} Memrise: официальный сайт платформы для изучения языков. — URL: https://www.memrise.com (дата обращения: 16.09.2025). — Текст : электронный.

\bibitem{} Puzzle English: официальный сайт сервиса для изучения английского языка. — URL: https://puzzle-english.com (дата обращения: 18.09.2025). — Текст : электронный.

\bibitem{} Quizlet: официальный сайт сервиса учебных карточек. — URL: https://quizlet.com (дата обращения: 18.09.2025). — Текст : электронный.

\bibitem{} Android Developers: официальная документация по платформе Android. — URL: https://developer.android.com (дата обращения: 05.10.2025). — Текст : электронный.

\end{thebibliography}

\chapter*{Приложение}
\addcontentsline{toc}{chapter}{Приложение}
\markboth{Приложение}{Приложение}

\subsection*{Приложение 1}

Фрагмент лексического набора по теме «Деловая коммуникация»

\begin{center}
\small
\begin{tabular}{|p{2.6cm}|p{2.9cm}|p{3.0cm}|p{7.6cm}|}
\hline
Слово & Транскрипция & Перевод & Пример \\
\hline
deadline & {\ipa [ˈdedlaɪn]} & крайний срок & The deadline for the report is Friday. \\
\hline
invoice & {\ipa [ˈɪnvɔɪs]} & счёт на оплату & Please send the invoice by e-mail. \\
\hline
negotiate & {\ipa [nɪˈɡəʊʃieɪt]} & вести переговоры & We need to negotiate the price. \\
\hline
agenda & {\ipa [əˈdʒendə]} & повестка дня & The agenda has five items. \\
\hline
colleague & {\ipa [ˈkɒliːɡ]} & коллега & My colleague will join the meeting. \\
\hline
schedule & {\ipa [ˈʃedjuːl]} & график, расписание & The schedule changed again. \\
\hline
estimate & {\ipa [ˈestɪmət]} & смета, оценка & The estimate looks too high. \\
\hline
profit & {\ipa [ˈprɒfɪt]} & прибыль & The company reported a small profit. \\
\hline
supplier & {\ipa [səˈplaɪə]} & поставщик & The supplier missed the delivery date. \\
\hline
confirm & {\ipa [kənˈfɜːm]} & подтверждать & Please confirm the time of the call. \\
\hline
attach & {\ipa [əˈtætʃ]} & прикреплять & I attach the signed contract. \\
\hline
postpone & {\ipa [pəˈspəʊn]} & переносить & We had to postpone the meeting. \\
\hline
\end{tabular}
\end{center}

\subsection*{Приложение 2}

Примеры упражнений, формируемых приложением автоматически

\begin{center}
\small
\begin{tabular}{|p{4.2cm}|p{12.2cm}|}
\hline
Тип упражнения & Пример задания \\
\hline
Выбор перевода & deadline — 1) срок сдачи; 2) задержка; 3) доставка; 4) договор \\
\hline
Выбор слова по переводу & «счёт на оплату» — 1) invoice; 2) receipt; 3) bill; 4) check \\
\hline
Ввод с клавиатуры & Подтверждать — c \_ \_ f \_ \_ m (подсказка: первая буква) \\
\hline
Пропуск в предложении & We had to \_\_\_\_\_\_ the meeting until Monday. (postpone) \\
\hline
Сборка предложения & the / invoice / send / please / by / e-mail \\
\hline
Перевод фразы & Пожалуйста, подтвердите время звонка. — Please confirm the time of the call. \\
\hline
\end{tabular}
\end{center}

\subsection*{Приложение 3}

Анкета участника экспериментальной группы

\begin{enumerate}

  \item Сколько раз в неделю вы открывали приложение? (менее двух, два-три, четыре и более)

  \item Удобно ли вам работать с приложением в транспорте и на перемене?

  \item Стали ли вы заниматься языком регулярнее, чем до начала работы с приложением?

  \item Какие упражнения показались наиболее полезными? (выбор перевода, ввод слова, пропуск в предложении, сборка предложения, перевод фразы)

  \item Что в работе приложения раздражало или мешало?

  \item Хотели бы вы продолжать пользоваться приложением в следующем семестре?

\end{enumerate}

\end{document}
